% Part: sets-functions-relations
% Chapter: sets
% Section: basics

\documentclass[../../../include/open-logic-section]{subfiles}

\begin{document}

\olfileid[tr]{sfr}{set}{bas}
\olsection{Kaplamsallık}

Bir \emph{küme}, tek bir nesne olarak ele alınan nesneler topluluğudur.
Kümeyi oluşturan nesnelere kümenin \emph{elemanları} veya \emph{üyeleri}
denir. Eğer $x$, $A$ kümesinin elemanıysa $x \in A$ yazarız; değilse
$x \notin A$ yazarız. Hiç elemanı olmayan kümeye \emph{boş küme} denir ve
bu küme ``$\emptyset$'' ile gösterilir.

\begin{explain}
Kümeyi nasıl \emph{belirttiğimizin}, elemanlarını hangi \emph{sırayla}
yazdığımızın, hatta aynı elemanı \emph{kaç kez} yazdığımızın önemi yoktur.
Önemli olan yalnızca elemanlarının hangileri olduğudur. Bunu aşağıdaki ilkeyle
kesinleştiririz.
\end{explain}

\begin{defn}[Kaplamsallık]
  $A$ ve $B$ kümelerse, $A = B$ eşitliği ancak ve ancak $A$ kümesinin her
  elemanı aynı zamanda $B$ kümesinin de elemanıysa ve bunun tersi de
  geçerliyse sağlanır.
\end{defn}

Kaplamsallık bazı gösterimleri kullanmamıza olanak verir. Genel olarak
$a_{1}$, \dots, $a_{n}$ nesneleri verildiğinde $\{a_{1}, \dots, a_{n}\}$,
elemanları $a_1, \ldots, a_n$ olan \emph{tek} kümedir. ``\emph{Tek}''
sözcüğünü vurguluyoruz; çünkü kaplamsallık böyle bir kümeden yalnızca
\emph{bir} tane olabileceğini söyler. Gerçekten kaplamsallık aşağıdakine de
izin verir:
  \[
    \{a, a, b\} = \{a, b\} = \{b,a\}.
  \]
Bu eşitlik, kümeler söz konusu olduğunda elemanların sırasını veya kaç kez
belirtildiklerini önemsemediğimizi gösterir.

\begin{tagblock}{novice}
\begin{ex}
Elinizde bir grup nesne olduğunda onları bir kümede toplayabilirsiniz.
Örneğin Richard'ın kardeşleri bir kişiyi içeren bir küme oluşturur ve bunu
$S=\{\textrm{Ruth}\}$ diye yazabiliriz. $4$'ten küçük pozitif tam sayıların
kümesi $\{1, 2, 3\}$'tür; ancak $\{3, 2, 1\}$, hatta $\{1, 2, 1, 2, 3\}$
biçiminde de yazılabilir. Kaplamsallık gereğince bunların tümü aynı kümedir.
Çünkü $\{1, 2, 3\}$ kümesinin her elemanı, $\{3, 2, 1\}$ kümesinin de
(ve $\{1, 2, 1, 2, 3\}$ kümesinin de) elemanıdır; bunun tersi de geçerlidir.
\end{ex}
\end{tagblock}

Bir kümeyi çoğu kez elemanlarının ortak bir özelliğiyle
belirteceğiz. Bunun için şu kısa gösterimi kullanacağız:
$\Setabs{x}{\phi(x)}$. Burada $\phi(x)$, $x$ nesnesinin kümenin elemanları
arasında sayılabilmesi için sahip olması gereken özelliği gösterir.

\begin{tagblock}{novice}
\begin{ex}
Örneğimizde $S$ kümesini şöyle de belirtebilirdik:
\[
S = \Setabs{x}{x \text{ Richard'ın kardeşidir}}.
\]
\end{ex}
\end{tagblock}

\begin{tagblock}{math}
\begin{ex}
Bir sayı, öz bölenlerinin (yani sayıyı kalansız bölen ama sayının kendisi
olmayan sayıların) toplamına eşitse \emph{mükemmel} olarak adlandırılır.
Örneğin $6$ mükemmeldir; çünkü öz bölenleri $1$, $2$ ve~$3$'tür ve
$6 = 1 + 2 + 3$ olur. Gerçekte $6$, $10$'dan küçük tek mükemmel pozitif tam
sayıdır. Dolayısıyla kaplamsallığı
kullanarak şöyle diyebiliriz:
\[
	\{6\} = \Setabs{x}{x\text{ mükemmeldir ve }0 \leq x \leq 10}
\]
Sağdaki gösterimi, ``$x$ mükemmeldir ve $0 \leq x \leq 10$ koşulunu sağlayan
tüm $x$ değerlerinin kümesi'' diye okuruz. Buradaki özdeşlik, kümeler söz
konusu olduğunda onların nasıl belirtildiğini önemsemediğimizi doğrular. Daha
genel olarak kaplamsallık, $\phi(x)$ özelliğini taşıyan $x$ değerlerinin
kümesinin her zaman yalnızca bir tane olduğunu güvence altına alır. Böylece
kaplamsallık, $\Setabs{x}{\phi(x)}$ ifadesini $\phi(x)$ özelliğini taşıyan
$x$ değerlerinin \emph{kümesi} diye adlandırmamızı haklı çıkarır.
\end{ex}
\end{tagblock}

Kaplamsallık, kümelerin özdeş olduğunu göstermek için bize bir yol verir:
$A = B$ olduğunu göstermek için, $x \in A$ olduğunda $x \in B$ olduğunu ve
$y \in B$ olduğunda $y \in A$ olduğunu gösterin.

\begin{prob}
En fazla bir boş küme bulunduğunu ispatlayın; başka bir deyişle, $A$ ve $B$
hiç elemanı olmayan kümelerse $A = B$ olduğunu gösterin.
\end{prob}

\end{document}
